\subsection{World Models for Reinforcement Learning and Embodied AI}

The idea of learning an internal model of environment dynamics dates back to early work on mental simulation and forward models in cognitive science and control theory. In deep reinforcement learning, \citet{ha2018worldmodels} proposed learning a compressed spatial and temporal representation of the environment using a variational autoencoder and an RNN, training a compact policy entirely within the learned ``dream.'' PlaNet \citep{hafner2019planet} formalised this via a latent dynamics model (RSSM) that plans directly in latent space from pixel observations. SimPLe \citep{kaiser2020simple} demonstrated model-based Atari play by training a video-prediction model as a learned simulator. MuZero \citep{schrittwieser2020muzero} showed that a learned dynamics model with Monte-Carlo tree search can master board games and Atari without access to the ground-truth rules.

The Dreamer family \citep{hafner2020dreamerv1, hafner2021dreamerv2, hafner2025dreamerv3} progressively refined the RSSM-based world model, culminating in DreamerV3 \citep{hafner2025dreamerv3}. They achieve human-level performance across over 150 diverse tasks with a single set of hyperparameters. Transformer-based world models subsequently emerged: IRIS \citep{micheli2023iris} replaced the recurrent backbone with an autoregressive transformer over discrete tokens; TransDreamer \citep{chen2022transdreamer} introduced a Transformer State-Space Model for memory-demanding tasks; $\Delta$-IRIS \citep{micheli2024deltairis} improved tokenization efficiency via context-aware delta encoding; DIAMOND \citep{alonso2024diamond} leveraged diffusion models to produce visually faithful world simulations; and EMERALD \citep{burchi2025emerald} achieved state-of-the-art Crafter performance using masked generative transformers over spatial latent states.

At a larger scale, video generation models have been cast as world simulators. OpenAI's Sora \citep{openai2024sora} demonstrated long-form video generation with emergent 3D consistency, while Genie \citep{bruce2024genie} and Genie~3 \citep{deepmind2025genie3} showed that text-conditioned generative models can produce interactive, explorable environments. Several surveys chart the rapid expansion of world models into autonomous driving \citep{wang2025wmadsurvey, wang2024wminitialsurvey}, embodied AI \citep{zhu2025embodiedwmsurvey}, and video generation \citep{cho2024sorasurvey, wang2025mechanisticvideo}.

A persistent challenge across all these approaches is \emph{compounding prediction error}: small inaccuracies at each rollout step accumulate exponentially over long horizons, degrading trajectory fidelity \citep{xiao2019compoundingerror, talvitie2017self, lambert2022survey}. Various mitigation strategies have been proposed, including short-horizon re-planning, self-correcting models \citep{talvitie2017self}, and physics-informed architectures \citep{pinwm2025, prophy2025}, yet the fundamental tension between computational depth and rollout stability remains unresolved by existing architectures.


\subsection{Looped and Recurrent-Depth Transformer Architectures}

Looped transformers reuse a shared set of transformer blocks across depth, decoupling effective computation from parameter count. The Universal Transformer \citep{dehghani2019universal} first proposed this idea, combining weight sharing with Adaptive Computation Time (ACT) \citep{graves2016act} for input-dependent halting. ALBERT \citep{lan2020albert} demonstrated the practical viability of full cross-layer parameter sharing in BERT-scale models.

Theoretical analyses subsequently established the computational power of looped transformers. \citet{giannou2023programmable} proved that looped transformers can simulate arbitrary programs, functioning as programmable computers with constant parameter count.  \citet{yang2023looped} showed that looped transformers match standard transformer performance on in-context learning while using less than 10\% of the parameters. \citet{fan2024lengthgen} demonstrated significant length generalisation improvements through adaptive loop counts. \citet{saunshi2025latentthoughts} provided both theoretical and empirical evidence that looped models implicitly generate ``latent thoughts,'' enabling reasoning beyond their apparent depth. At a practical scale, Ouro \citep{zhu2025ouro} trained looped language models (LoopLMs) through the full modern LLM pipeline with pre-training, instruction tuning, and RLHF, achieving 2--3$\times$ parameter efficiency with entropy-regularised adaptive computation. \citet{geiping2025rdm} demonstrated that recurrent-depth models (RDMs) scale test-time compute by increasing loop count at inference, following predictable quality improvements. \citet{pappone2025twoscale} analysed the geometry of latent dynamics in recurrent-depth transformers, identifying two-scale structure with fast intra-loop and slow inter-token dynamics. LoopFormer \citep{jeddi2025loopformer} introduced elastic-depth training with shortcut modulation for budget-conditioned reasoning. Mixture-of-Recursions \citep{bae2025mor} proposed per-token dynamic depth allocation within a single recursive framework. MoEUT \citep{csordas2024moeut} combined mixture-of-experts with universal transformers to balance specialisation and sharing.

\subsection{Adaptive Computation and Early Exit}

Allocating variable computation to inputs of differing complexity has been studied across multiple paradigms. \citet{graves2016act} introduced Adaptive Computation Time for RNNs, allowing per-step halting decisions. The early exit literature \citep{teerapittayanon2016branchynet, bolukbasi2017adaptive, earlyexit2025survey} enables inference to terminate at intermediate layers when confidence is sufficient. In the context of looped transformers, adaptive depth takes a particularly natural form: the model can halt after any number of loop iterations. Ouro \citep{zhu2025ouro} introduced entropy-regularised early exit, where a token exits the loop when its prediction entropy drops below a learned threshold. \citet{geiping2025rdm} trained with stochastic depth sampling (Poisson-distributed loop counts) to induce robustness to variable test-time depth. LoopFormer \citep{jeddi2025loopformer} conditioned on a continuous ``time budget'' during training, enabling fine-grained compute allocation at inference. \citet{pappone2025twoscale} proposed acceleration-based exit rules using second-order differences of hidden states. For world models specifically, adaptive computation is highly attractive: simple state transitions (e.g., free flight, static scenes) demand minimal processing, while complex events (e.g., multi-body collisions, contact dynamics) require deeper iterative refinement. To the best of our knowledge, no prior work has proposed adaptive-depth looped architectures with world modelling.


